\SourceLeaf{31} time interval $\mathrm{d}t$ as a result of the forces acting on them during that time and the speeds and directions they had attained at time $t$, if they were all completely free. The actual positions $c,c',c''$ and so on will then be those, among all positions compatible with the conditions of the system, for which $m(bc)^2+m'(b'c')^2+m''(b''c'')^2$ and so on becomes a minimum.

Equilibrium is evidently only a particular case of the general law. Its condition is that
\[
m(ab)^2+m'(a'b')^2+m''(a''b'')^2\ \text{and so on}
\]
itself be a minimum, or that the system’s remaining at rest come closer to the free motion of the individual points than any possible departure from that state.

Our principle is readily derived from the principles cited above in the following way.

The force acting on the material point $m$ is evidently made up of two forces. The first, together with the speed and direction at time $t$, takes the point in the time $\mathrm{d}t$ from $a$ to $c$. The second would, in the same time, take it from rest at $c$ along $cb$ if the point were regarded as free. The same holds for the other points. By \textsc{d’Alembert}’s principle, the points $m,m',m''$ and so on must therefore be in equilibrium under the action of the second forces alone, directed along $cb,c'b',c''b''$ and so on, at the positions $c,c',c''$ and so on, by virtue of the conditions of the system.

\expandafter\def\csname HistoricalNoteText3\endcsname{The usual formulation tacitly assumes conditions under which the reverse of every possible motion is also possible. A point may, for instance, be required to stay on a given surface, or the distance between two points may be fixed, and so forth. But this restriction is unnecessary and does not always fit nature. The surface of an impenetrable body does not compel a material point on it to stay there. It merely prevents the point from passing to one side. A taut thread between two points, inextensible but flexible, prevents only an increase in their distance and allows a decrease, and so on. Why then should we not state the law of virtual velocities from the outset so that it covers \textit{all} cases?}
By the principle of virtual velocities, this equilibrium requires that the sum of the products of three factors for each point, namely the corresponding masses $m,m',m''$ and so on, the line segments $cb,c'b'\ c''b''$ and so on, and any motions of those points that the conditions of the system allow, projected onto the corresponding line segments, always be $=0$, as it is usually put\HistoricalNote{3}
